%% ============================================================
%% SECTION 10: CONCLUSION
%% EarthVision-LM Survey -- Artificial Intelligence Review
%% Template: sn-jnl (Springer Nature) | sn-mathphys-num
%% ============================================================

\section{Conclusion}\label{sec10}

This survey has provided a comprehensive, systematic, and architecturally
grounded analysis of multimodal large language models for remote sensing.
We propose an architecture-first five-dimensional design space for synthesising
RS-MLLMs across encoder adaptation, connector design, LLM strategy, modality scope,
and spatial output---a framework that, to the best of our knowledge, no existing
RS-MLLM survey has previously formalised with explicit mapping of RS domain
challenges to each design dimension. The analysis covers
97 primary RS-MLLM papers (Corpus~A) supported by 113 contextual
domain references (210 total retrieved), published between January 2022 and August 2026. We organised the
analysis around a central diagnostic question: \emph{why do RS-MLLMs in 2026 remain
perception-limited despite three years of rapid development?} The answer, as
revealed by the five-dimensional design-space analysis, multi-modal coverage study,
hallucination taxonomy, and quantitative synthesis, is structural rather than incidental.

\textbf{The structural findings.}
Three structural gaps, each traceable to a specific dimension of the RS-MLLM design
space, account for the perception gap between current systems and operational
geospatial intelligence.

The \emph{modality gap} ($\mathcal{M}$-dimension) is the most severe: 84.6\,\% of
published RS-MLLMs process optical imagery only, despite RS being inherently
multi-sensor. Recent models such as EarthGPT~\cite{earthgpt2025},
EarthGPT-X~\cite{earthgptx2026}, and Earth-OneVision~\cite{earthonevision2026}
substantially expand sensor coverage---Earth-OneVision reaching optical, SAR,
infrared, multispectral, temporal, and video across nine task categories.
However, broad HSI capability, physics-grounded cross-sensor reasoning, and
robust spatial reasoning remain underdeveloped across the field.
The four specific structural deficits are:
(i)~broad HSI coverage (no dedicated HSI RS-MLLM exists);
(ii)~physics-specific modality encoding for SAR backscatter and hyperspectral
spectral signatures;
(iii)~systematic cross-sensor evaluation protocols; and
(iv)~operationally grounded spatiotemporal reasoning.
HM-Bench results are consistent with the hypothesis that RGB-oriented
representations and inadequate spectral grounding contribute to Type~III failures,
confirmed across all 18 evaluated models. These gaps cannot be closed by scaling
optical training
data---they require modality-specific encoder design and physics-informed training
objectives.

The \emph{spatial output gap} ($\mathcal{S}$-dimension) is the most operationally
consequential: oriented bounding box (OBB) output---the required representation for
aircraft, ships, and vehicles at arbitrary rotation angles---is present in four
Corpus~B models (GeoChat, EarthGPT, GeoGround, EarthGPT-X) but limited to optical
imagery in all cases. Cross-sensor OBB reasoning (SAR/HSI OBB) is essentially
unaddressed across the entire Corpus~B. Level-5 cross-sensor OBB with segmentation
is entirely unoccupied.

The \emph{task gap} (Level-4 spatiotemporal reasoning) is the most aspirational:
operational RS requires multi-temporal causal reasoning; while Earth-OneVision~\cite{earthonevision2026}
supports temporal inputs and VLRS-Bench~\cite{vlrsbench2026} now targets
complex reasoning across up to eight temporal phases, explicit causal and
longitudinal reasoning remains substantially underdeveloped, despite these
emerging temporal and multi-phase benchmarks and model capabilities. The field
has saturated Level-1 image-level tasks (VQA, captioning) but remains nascent
at Level-3 pixel-level tasks and substantially underdeveloped at Level-4.

\textbf{The path forward.}
The eight future research directions charted in Sect.~\ref{sec9} provide a concrete
roadmap from perception-limited systems to geospatially intelligent ones. The
near-term directions---OBB-capable MLLM training, a Type~III hallucination
benchmark, and a physics-informed SAR encoder---are tractable within the current
architectural paradigm. Relevant prior work on long-tail OBB
detection~\cite{odfe2025}, physics-constrained SAR--multispectral
fusion~\cite{rayleigh2026}, and semi-supervised temporal consistency~\cite{sm2lf2026}
provides methodological foundations for these directions.

The mid-term directions require architectural innovation: unified multi-sensor
MLLMs with modality-specific encoder banks, spectral-physics-aware HSI encoders,
and spatiotemporal reasoning with temporal token fusion and change-aware attention.
The long-term vision---a geospatial foundation model pre-trained across all sensor
modalities, spatial resolutions, and temporal scales---requires compute and data
infrastructure that is currently beyond academic budgets but represents the natural
convergence point of the trends documented in this survey.

\textbf{What this survey contributes to the field.}
The five-dimensional design space (Sect.~\ref{sec3}) provides the RS-MLLM community
with a shared vocabulary for architectural comparison and gap identification. The
task hierarchy (Sect.~\ref{sec4}) provides a framework for prioritising evaluation
investment. The multi-modal coverage analysis (Sect.~\ref{sec5}) quantifies the
modality debt---the investment in SAR and HSI infrastructure that the field must
make to realise the full potential of RS-MLLM technology. The three-type
hallucination taxonomy (Sect.~\ref{sec7}) provides the conceptual foundation for
hallucination detection and mitigation across all sensor modalities. The structured quantitative synthesis
(Sect.~\ref{sec8}), grounded in a PRISMA-compliant search yielding 97 primary
RS-MLLM papers (Corpus~A: 31 model studies yielding 26 architectures; 18 benchmark
studies; 14 dataset studies; 9 hallucination studies; 25 foundation-model studies)
supported by 113 contextual domain references, establishes
the empirical baseline against which future field growth can be measured.
These contributions are explicitly complementary to the concurrent survey of
Ma~et~al.~\cite{ma2026dsorsp}, which comprehensively covers RS-MLLM evolution,
multimodal learning, training data, downstream tasks, benchmark comparison,
generalisation, and spatial/UHR reasoning, with a primary contribution of
empirical RS-vs-CV-MLLM diagnostic evaluation demonstrating that general-purpose
CV-MLLMs can match RS-specific models on several tasks. Ma~et~al.\ establish
\emph{what} the performance gaps are empirically; EarthVision-LM explains
\emph{why} they exist at the architectural and physics level, and charts the
path to closing them. Readers are directed to both surveys for a complete
picture of the RS-MLLM landscape.

\textbf{The broader significance.}
Remote sensing is among the most consequential application domains for AI,
spanning high-consequence applications including disaster response, infrastructure
monitoring, agricultural management, and environmental surveillance.
RS-MLLMs that hallucinate confidently are not merely inaccurate systems:
they are potentially dangerous systems, as established in Sect.~\ref{sec7:danger}.
The development of RS-MLLMs that are physically grounded, multi-sensor capable,
spatially precise, and hallucination-aware is therefore not only a scientific
challenge but a societal one.

The field is young---the oldest dedicated RS-MLLM (RSGPT~\cite{rsgpt2023}) appeared
less than three years before this survey was written---and its growth trajectory
(3 papers in 2022, 42 in 2025) suggests that the structural gaps identified here
will be actively contested in the coming years. We hope that the taxonomy, analysis,
and roadmap presented in EarthVision-LM will provide a useful compass for that work.
